\section{Introduction}
A renewal provider may distinguish many service histories while its terminal gateway accepts only a small menu of commands. The provider must choose that menu and allocate an accepted obligation across histories. Participation caps restrict expected service, realization ceilings restrict every positive-probability delivery, and opening a command may incur a fixed charge. The menu and the service targets are therefore joint decisions. A fast menu optimizer for prescribed targets does not by itself solve the operational problem.

We show that heterogeneous histories need not create heterogeneous search coordinates. Histories that admit the same catalog commands can be pooled exactly, even when every participation cap and shortfall cost differs. The reason is a terminal-first allocation rule. Within an eligibility class, intermediate service and unfinished terminal capacity cannot coexist at an optimum. Before terminal capacity is exhausted, capped equalization allocates the group's obligation. Afterward, a separable capped-cost problem allocates the remaining service. Its dependence on the menu is confined to the last eligible command.

This localization changes the joint optimization problem. We add one exact correction at the path edge that leaves an eligibility class, or at the final selected command. A conditional path algorithm then optimizes both the menu and all individual targets for supplied class means. With common catalog eligibility, the accepted root obligation already fixes the sole class mean. One polynomial-time call solves the original finite-catalog problem globally, at interior as well as saturated obligations, for any number of histories, variable menu size, heterogeneous caps and quadratic costs, and arbitrary nonnegative opening charges. No averaging error, cap guard, command augmentation, or target grid is required in this regime.

For several eligibility classes, we derive an exact price oracle on boxes of class means. The oracle combines the terminal-first correction with an increasing and a decreasing path through the common priced command scores. The second side is needed because mandatory lower targets can make commands after the price-optimal peak indispensable. Every returned upper bound covers all allowed menus in its box. Subdivision gives an original-space interval at every interruption and reaches any positive additive tolerance when completed. Its worst-case accuracy exponent depends on the number of distinct eligible prefixes, not on the number of histories or the number of cost bins. Dependence on eligibility complexity remains exponential; the common-eligibility regime is exactly polynomial and has no accuracy parameter.

The certificate is independently checkable. A feasible service allocation and a scalar price certify each pooling correction. Separate checks cover the path inequalities, eligible anchors, every branch of the target-space partition, and the original lotteries. The checker does not invoke the optimization algorithm or infer global accuracy from a solver status. Thus the large-history result is a joint optimization statement, not a conditional solution accompanied only by a Jensen bound.

Our computational study first revisits every instance in the preceding cost-heterogeneity experiment without changing its primitives. It then tests common eligibility with up to 1,024 distinct histories and stresses growing eligibility complexity at two accuracies and two node budgets. We retain all unresolved intervals, exact exhaustive references where specified, and every direct mixed-integer envelope solve. The experiment protocol precedes execution in the repository. These are declared synthetic inputs; no calibration to an operational population or implemented deployment is asserted. Historical solver times are distinguished from new measurements rather than presented as matched-machine speed comparisons.

\paragraph{Relation to established methods.}\label{sec:r37-position}
Interval quantization, matrix searching, pricing, and branch-and-bound supply important building blocks \citep{Wu1991,NielsenNock2014,GronlundEtAl2017,HandlerZang1980,LawlerWood1966,Lemarechal2001,HubnerEtAl2025}. Capped separable resource allocation and its multiplier algorithms are also established \citep{PatrikssonStromberg2015,SchootUiterkamp2022}. Our contribution is not another water-level solver or another general subdivision rule. It is the terminal-first identity and its last-eligible-command correction, which make these tools jointly optimize a shared charged alphabet without sacrificing individual constraints. The two-sided group-box identity supplies a locally exact upper oracle for the remaining eligibility means.

The earlier aggregation results address a complementary construction. Value-function aggregation and error decomposition have broad precedents \citep{LiBertsekas2025,Fattahi2025}, and arithmetic--harmonic mean inequalities are classical \citep{LiaoWu2015}. Our retained guarded and harmonic relaxations preserve the installed menu and allow an exact original-feasible lift. Exact eligibility pooling goes further by retaining the full within-class service frontier instead of replacing it by a representative curvature. It removes the resulting dispersion allowance altogether, at the cost of evaluating the original histories within each oracle call.

Canonical contractual allocation, continuous quadratic design, Monge acceleration, budgeted compression, target-net approximation, and charge-aware recovery remain part of the mathematical contribution under their original assumptions. Table~\ref{tab:r49-complexity} distinguishes these regimes from the exact eligibility results. The main text develops the structural joint algorithm; regular appendices and the electronic companion retain the earlier proofs. A separate computational record preserves all earlier study narratives and supplies the complete per-instance measurements and reproducibility instructions.
